\begin{table}[htbp]
\centering
\small
\caption{Leave-one-out validation on unseen episodes. Each episode was held out in turn and the model built from all the others. Means over held-out episodes; the difference is second-order minus first-order on the same episode, with its 95\% CI.}
\label{tab:loo}
\setlength{\tabcolsep}{3pt}
\begin{tabular*}{\linewidth}{@{\extracolsep{\fill}}lrrrr}
\toprule
 & \multicolumn{3}{c}{\textbf{Order}} & \\
\cmidrule(lr){2-4}
\textbf{Score} & \textbf{0} & \textbf{1} & \textbf{2} & \textbf{Difference (2 $-$ 1) [95\% CI]} \\
\midrule
\multicolumn{5}{l}{\textit{Predicting the next contact}} \\
Log-likelihood (5 cat.) & -1.430 & -1.123 & -1.087 & +0.036 [+0.008, +0.064] \\
Perplexity (5 cat.) & 4.18 & 3.07 & 2.97 & \\
Log-likelihood (7 cat.) & -1.697 & -1.357 & -1.312 & +0.046 [+0.022, +0.070] \\
Perplexity (7 cat.) & 5.46 & 3.89 & 3.71 & \\
Accuracy (5 cat.) & 0.446 & 0.592 & 0.637 & +0.044 [+0.014, +0.075] \\
Accuracy (7 cat.) & 0.396 & 0.520 & 0.556 & +0.037 [+0.013, +0.060] \\
\addlinespace
 & \textbf{Scrambled} & & & \\
\multicolumn{5}{l}{\textit{Similarity of simulated chains to the held-out episode}} \\
Fixed, Lev.\ ($\downarrow$) & 0.610 & 0.615 & 0.631 & +0.016 [+0.014, +0.019] \\
Fixed, S-W ($\uparrow$) & 0.344 & 0.371 & 0.362 & -0.009 [-0.012, -0.006] \\
Fitted, Lev.\ ($\downarrow$) & 0.650 & 0.654 & 0.666 & +0.012 [+0.009, +0.015] \\
Fitted, S-W ($\uparrow$) & 0.379 & 0.408 & 0.399 & -0.009 [-0.012, -0.005] \\
Terminating, Lev.\ ($\downarrow$) & 0.671 & 0.620 & 0.609 & -0.011 [-0.014, -0.007] \\
Terminating, S-W ($\uparrow$) & 0.384 & 0.440 & 0.422 & -0.018 [-0.022, -0.014] \\
\bottomrule
\end{tabular*}
\par\smallskip
\begin{minipage}{\linewidth}\footnotesize
Log-likelihood is per transition (higher is better) and perplexity $= \exp(-$log-likelihood$)$ (lower is better). Accuracy is the share of next contacts predicted correctly, with shared credit for ties. Only transitions with at least two earlier contacts are scored, so both orders are scored on the same transitions; 154 and 163 episodes had at least one such transition with five and seven categories. Transition probabilities use add-one smoothing; with smoothing 0.1 the five-category log-likelihood difference is +0.019 [-0.012, +0.049]. Zero-order predicts from category frequencies only. For the similarity rows the first column is the scrambled baseline (the first-order chains with their contacts shuffled) instead of the zero-order model. Similarity uses 100 simulated chains per held-out episode; Lev.\ = normalised Levenshtein distance (lower is better), S-W = normalised Smith-Waterman score (higher is better).
\end{minipage}
\end{table}
